\scanpage{18}

à des $f$ de type fini au lieu de $f$
propre. Se ramène au cas d'une \uline{imm.}
\uline{ouverte}. Se résoud par l'affirmative
lorsqu'on dispose de la résolution des
singularités \inserted{$+$ pureté}, p.ex. au cas \unsure{des}
\inserted{pour ``bons préschémas''}
\unsure{espaces} \inserted{\unsure{bons présch}}
alg. ou lorsqu'on dispose de résolution
des singularités \deleted{\unsure{de}} dans ce cadre\ldots

14) \uline{Les} $R^i_!f_*(F)$.

Soit $f:X\longrightarrow Y$ morphisme compactifiable,
i.e. qui se factorise en
\diagram{\begin{tikzcd}
X\arrow[r,"i","{\mathrm{imm.\ ouv.}}"']&
\overline X\arrow[r,"\bar f","{\mathrm{propre}}"']&Y
\end{tikzcd}}
(p.ex. $f$ qu. proj. et $Y$ quasi-cpt).

Pour $\forall$ faisceau de torsion $F$ sur $X$,
considérons $i_!(F)$ [ex. et en $F$]
puis les $R^q_{\deleted{\unsure{!}}}\bar f_*(i_!(F))$, forment
un foncteur coh. en $F$. Il résulte
de façon quasi-triviale du th. du changement
de base que ce foncteur coh.

\mspage{17}
